\section{Introduction}
Device pass-through
has emerged as an essential architectural mechanism 
for delivering near-native I/O performance 
to Virtual Machines (VMs).
It can be realized through direct device assignment,
Single Root I/O Virtualization (SR-IOV)~\cite{ms_sriov_overview, intel_sriov_ethernet, amd_sriov_mxgpu}
to the increasingly granular Scalable I/O Virtualization (Scalable IOV)~\cite{intel_scalable_iov_spec}.
In such virtualized environments, 
the physical Input/Output Memory Management Unit (IOMMU) 
plays a critical role in enforcing strong isolation between guests. 
It effectively prevents a hardware device---or a specific virtual function---assigned to one VM 
from performing unauthorized Direct Memory Access (DMA) 
to the memory of another.

Despite these vital inter-guest safeguards, 
traditional pass-through mechanisms 
suffer from a lack of \emph{intra-guest protection}\cite{Amit2011vIOMMU,Tian2020coIOMMU,Lv2022LAvIOMMU}. 
Specifically, an assigned device is granted 
uncontrolled access to the entire guest physical address space. 
This broad access paradigm leaves the guest operating system 
highly vulnerable to memory corruption or data exfiltration, 
whether originating from buggy device drivers 
or fundamentally malicious hardware components\cite{Amit2011vIOMMU, RajapakshaIOMMU, AlexCodeInjection}.

Furthermore, from a hypervisor resource management perspective, 
device pass-through conventionally necessitates 
the static pinning of all guest memory pages. 
This rigid architectural requirement 
inherently precludes memory overcommitment \cite{Wang2025ToPRI,Guo2024VPRI}, 
thereby severely degrading overall system resource utilization.

To overcome these intra-guest security vulnerabilities 
and hypervisor management limitations, 
a Virtualized IOMMU (vIOMMU) 
has been proposed and studied in both literature \cite{Amit2011vIOMMU,Tian2020coIOMMU} and practice \cite{qemu_10_0,qemu_vtd}. 
vIOMMU exposes a virtual IOMMU to the guest OS so that unmodified guest OS can defend itself against buggy device drivers and malicious hardware.

However, existing vIOMMU designs
have been reported with severe performance degradation~\cite{Wang2025ToPRI,Tian2020coIOMMU,Amit2011vIOMMU}.
% suffer from severe performance degradation~\cite{Wang2025ToPRI,Tian2020coIOMMU}.
Our profiling reveals that running vIOMMU may suffer from up to 90\% throughput degradation compared with no vIOMMU.

\textbf{The vIOMMU Bottleneck Pivot.}
\schaic{@Saksham, I rephrase your statement saying that the translation overhead is the bottleneck in baremetal.
Both F\&S and DAMN 
reduce both address translation and invalidation overhead. How can we claim that IOMMU address translation overhead was \emph{the} primary performance bottleneck? }
% In bare-metal environments, prior IOMMU studies demonstrated
% that IOMMU address translation overhead was the primary performance bottleneck 
%     when enforcing strict-level security.
In bare-metal environments, prior IOMMU studies demonstrated
that synchronous management overheads
---including both address translation and hardware IOTLB invalidation---were the primary performance bottlenecks.
This is especially true when enforcing strict-level security, the highest level of security in Linux kernel,
    which requires the immediate unmapping and IOTLB invalidation after each DMA completion.
% Strict-level security is the highest level of security in Linux kernel; it requires immediate invalidation of IOTLB entries upon unmapping.
Recent software optimizations have largely resolved this in bare-metal environments;
for instance, DAMN~\cite{markuze2018damn} uses software cache pools for DMAs,
and F\&S~\cite{Rubin2024FastSafe} batches mapping requests.

However, when migrating these strict-security models to virtualized environments,
a fundamentally new bottleneck emerges: unavoidable VM exits.
Early vIOMMU implementations relied on shadow page tables~\cite{qemu_vtd,Amit2011vIOMMU},
which suffered from VM exits on \emph{every} mapping and unmapping operation
to notify the hypervisor of page table changes.
To alleviate this, modern architectures introduced hardware-assisted nested IOMMU support~\cite{intel_vtd_spec},
which successfully eliminates VM exits during the creation of new IOMMU mappings.
Yet, to maintain strict-level security,
this hardware-assisted approach still mandates an IOTLB invalidation---and thus a synchronous VM exit---at unmap time~\cite{qemu_10_0, intel_vtd_spec_caching}.
Because the physical IOMMU invalidation queue is a shared host resource,
allowing multiple VMs to write to it directly would lead to corruption.

Consequently, the performance impact of these VM exits
completely dominates the system's I/O throughput.
A VM exit is inherently slow because it mandates a heavyweight hardware context switch
that pauses guest execution and saves its architectural state,
before trapping into the hypervisor to process the event in the host kernel.
In our profiling, an IOTLB invalidation that triggers a VM exit
takes between 10 to 20 microseconds.
This is orders of magnitude slower than a bare-metal hardware invalidation and address translation,
    which typically completes in 200 to 300 nanoseconds on modern hardware~\cite{Rubin2024FastSafe}.
\todo{more citations}

% The performance impact of these VM exits
% dominates the performance of the system.
% In our profiling, a VM exit triggered by IOTLB invalidation can take bewteen 10 - 20 microseconds.
% A VM exit is inherently slow because 
% it mandates a heavy-weight hardware context switch that pauses guest execution, 
% saves its architectural state, and traps into the hypervisor software stack to process the event.

% Intuitively, a single VM exit can consume thousands of CPU cycles,
% whereas a native memory operation takes only tens.
% Consequently, the per-unmap request latency increases drastically in a VM.
% Yet, despite this massive latency spike per request,
% the system's \emph{parallelism} remains artificially capped by legacy software synchronization,
% preventing the system from hiding this latency.
% Our profiling of the recent Linux kernel~\cite{linux_v6129}---which already adopts DAMN's page pool optimization---
% shows that even with hardware-assisted nested IOMMU,
% turning on vIOMMU still suffers up to 82\% throughput degradation
% on the Tx datapath compared with no vIOMMU.

% Recent architectures have developed dedicated support for IOMMU virtualization to address these performance challenges.
% Early vIOMMU implementations utilized shadow page tables, which suffered from high overhead due to the frequent VM exits required for IOTLB invalidation on every mapping and unmapping operation.
% These VM exits were necessary to notify the hypervisor and host kernel of each change\cite{Amit2011vIOMMU,qemu_vtd}.
% To improve performance, recent architectures introduced hardware-assisted Nested IOMMU support, 
%     which eliminates VM exits when the guest creates a new IOMMU mapping.
% However, this approach still requires IOTLB invalidation---and thus a VM exit---at unmap time~\cite{yiliu1765qemu}.
% This is because the IOMMU hardware is shared by all VMs, 
% and allowing multiple VMs to independently write to the invalidation queue would lead to corruption~\cite{intel_vtd_spec}.

\schaic{From previous version. Left for Tianyin's reference.\newline
The existing software design of hardware-assisted vIOMMU faces significant challenges in balancing security and performance.
Vanilla Linux kernel provides two invalidation policies: \emph{strict} and \emph{lazy}~\cite{linux_kernel_params}.
The strict invalidation policy provides maximum security by immediately purging IOTLB entries upon unmapping; however, it incurs a massive performance penalty.
% Our profiling shows that a guest VM running in strict mode achieves less than 15\% of the throughput compared to a system without vIOMMU enabled.
The lazy invalidation policy defers invalidation and batches requests within a deferred time window.
While more performant, this approach creates unignorable security risks regarding data confidentiality and integrity~\cite{RajapakshaIOMMU,AlexCodeInjection}.
}

% \schaic{This para is from previous version, left for Tianyin's reference.}
\schaic{From previous version. Left for Tianyin's reference.\newline
Other solutions attempt to optimize the software stack with driver changes but fail to close the performance gap with vIOMMU disabled.
For instance, DAMN~\cite{markuze2018damn} uses software cache pools for DMAs to reduce invalidations but is incompatible with zero-copy networking Tx paths~\cite{markuze2018damn}.
Similarly, F\&S~\cite{Rubin2024FastSafe} batches mapping requests but does not address the high cost of IOTLB invalidation in virtualized environments, 
    particularly on the Tx datapath where batching is more challenging.
Our profiling evaluates performance of recent version of Linux kernel~\cite{linux_v6129} which already adopts DAMN's page pool optimization.
The results show that even with hardware-assisted nested IOMMU,
    turning on vIOMMU still suffers up to 82\% throughput degradation compared with no vIOMMU.
}

\textbf{Key Insight.}
Our analysis shows that the existing software design fails to address the performance bottleneck
    given the expensive IOTLB invalidation in virtualized context. 
The current IOMMU subsystem design fully serializes all concurrent unmaps on the same domain,
    blocking future invalidation requests before the ongoing invalidatoin finishes.
    % single hardware round-trip finishes.
This serialization caps throughput independent of the core count.
Worse, scaling the number of cores actively degrades throughput since it triggers extra VM exits due to lock contentions.
This compounds hypervisor overhead and causes invalidation latency to grow super-linearly with the number of active cores.

Current software designs unnecessarily conflate
\emph{memory safety ordering} with \emph{cross-core serialization}.
To maintain strict protection guarantees,
it is only necessary to enforce a chronological ordering
for a \emph{specific} memory mapping within a single thread.
For a mapping from IO Virtual Address (IOVA) to guest physical address (GPA),
the IOVA must be unmapped, then invalidated in the hardware IOTLB, and only then the physical page freed and reused.
It does \emph{not} require serializing independent unmap requests across different cores.

% The root cause of performance degradation comes from the VM exits required for IOTLB invalidation on the Tx unmap path.
% Despite hardware-assisted nested translation eliminating VM exits on the mapping path,
%     strict-mode vIOMMU protection still exhibits severe performance degradation.

% With receive-side overheads largely mitigated by page pools,
%     the primary bottleneck shifts to transmit-path unmaps,
%     where IOTLB synchronization accounts for over 95\% of the unmap latency.
% Because the hardware invalidation queue is a shared host resource without per-VM partitioning,
%     every invalidation mandates a synchronous VM exit.


\textbf{Our approach.}
\schaic{@Saksham, I think carefully and avoid using the word "parallelism" here. 
"parallelism" implies multi cores handling invalidation requests in parallel. But we didn't have  }
To address these challenges, we rethink I/O memory protection mechanisms in VMs
to enable the system to process highly concurrent invalidation requests and, thus, amortize VM exit costs.
% To address these challenges, we rethink I/O memory protection mechanisms in VMs
% to vastly increase effective request parallelism.
We introduce \brand, which significantly improves vIOMMU performance
while maintaining security parity with the strict model.

First, we developed \emph{\optOne}, a technique that decouples
the submission of invalidation requests from the actual invalidation process.
A delegated core performs the IOTLB invalidation work,
while other cores submit invalidation requests and wait.
This decoupling removes the need to wait for previous invalidation to finish before generating an invalidation request,
destroying the cross-core serialization bottleneck.
Furthermore, it can process concurrent invalidation requests through efficient batching:
multiple invalidation requests can be coalesced
while the delegated core is processing a previous batch,
drastically reducing the total number of VM exits.
\emph{\optOne} requires zero changes to driver code;
all modifications are contained within the IOMMU subsystem.

% To address these challenges, we introduce \brand that 
% significantly improve vIOMMU performance while maintaining security parity with the strict model.
% First, we developed \emph{\optOne}, a technique that decouples the submission of invalidation requests from the actual invalidation process.
% A delegated core will perform the IOTLB invalidation work, while other cores submit invalidation requests
% and wait for their requests to finish.
% The decoupling of request submission and invalidation also reduces lock contention, 
%     as generating an invalidation request no longer requires holding the global lock.
% What's more, the design enables efficient batching of invalidation requests, 
%     as multiple invalidation requests can be coalesced together while the delegated core is processing a previous batch.
% The coalescence of invalidation requests reduce the number of VM exits, and thus significantly improves performance.
% \emph{\optOne} can be realized with zero change to any drive code; 
%     all changes are contained within the IOMMU subsystem, making it broadly applicable to all drivers.
%  offloads the invalidation process to a delegated core, 

% By decoupling the request submission from the invalidation process, we ensure that only the enqueue operation remains under . 
% By moving the IOTLB invalidation out of the locking protection, we significantly reduce lock contention, leaving only the enqueue operation under lock. 
% This design also enables \emph{Efficient Batching}, as multiple invalidation requests can be grouped together while the delegated core is processing a previous batch.

% However, to maintain the strict security model without sacrificing this newfound parallelism,
% the driver must not free the associated memory page
% until the asynchronous invalidation completes;
% otherwise, a Use-After-Free vulnerability arises~\cite{RajapakshaIOMMU}.
To maintain the strict security model, 
the driver must wait for the IOVA invalidation to complete \emph{before} freeing the associated guest physical address (GPA) page.
Otherwise, a the guest OS is vulnerable to data confidentiality attack\cite{RajapakshaIOMMU}.
To enforce strict safety ordering (Unmap $\rightarrow$ Invalidate $\rightarrow$ Free)
without forcing the CPU to synchronously wait,
we propose the \emph{Deferred Free Page (DFP)} mechanism.
DFP introduces a new callback-based interface in the IOMMU subsystem,
allowing drivers to pass a callback function
that delays physical memory cleanup until hardware invalidation is confirmed.
With DFP, the driver can proceed immediately without waiting for the invalidation to complete
while still preserving strict security guarantees.
% The cominbation of DFP and \optOne can enable the system to process more DMA packets since it spends less time in IOMMU subsystem.
% and, thus, improve the throughput without sacrificing strict security guarantees.
The cominbation of DFP and \optOne can enable the system to process more I/O requests and produce higher throughput
since the CPU spends significantly less time blocked in the IOMMU subsystem.

% However, simply delegating invalidation is not enough. 
% To maintain the strict security model, 
% the driver must wait for the invalidation to complete \emph{before} freeing the associated memory page to prevent data integrity and confidentiality issues.
% Drivers assume that a page that is logically unmapped is also physically unmapped, 
%     so the drivers may release the physical page before the hardware invalidation actually finishes.
% This can lead to security issues if the page is reused for a different purpose (e.g. store password) before the invalidation completes,
%     as the device could still access the page through DMA~\cite{RajapakshaIOMMU}.
% We remove the waiting time without lose of strict garrantee through
%     \emph{Deferred Free Page (DFP)} mechanism.
% Refactoring all drivers to handle this is impractical due to missing IOMMU subsystem interface and engineering scale. 
% To solve this, we propose the \emph{Deferred Free Page (DFP)} mechanism. 
% DFP introduces a new callback-based interface in the IOMMU subsystem.
% This allows drivers to pass a callback function and delay the physical memory cleanup until the hardware invalidation is confirmed;
%     as a result, the driver can proceed without waiting for the invalidation to complete, while still maintaining the strict security model.
% This technique is broadly applicable and compatible with both lazy and \optOne modes.
% With DFP combined with \optOne, we can achieve close-to-no-vIOMMU performance while maintaining the strict security model.

We implement \brand on top of Linux kernel 6.12.9
with fewer than 1.5k lines of code changed (only 373 lines in the device driver).
\brand achieves close-to-no-vIOMMU performance
while maintaining the strict security model,
yielding up to a 7x throughput improvement over the baseline strict model.

\textbf{Summary.} This paper makes the following contributions:
\begin{itemize}
    \item \textbf{Overhead characterization:} We provide a detailed analysis of the inefficiencies in existing vIOMMU designs,
    identifying cross-core serialization as the root cause of performance degradation under VM exit latency.
    \item \textbf{Close-to-no-vIOMMU performance:} We demonstrate that \brand achieves near-native performance
    by enabling the system to process highly concurrent invalidation requests and minimize blocking time in IOMMU subsystem.
    \item \textbf{Strict-level security:} We quantify the practical security compromises of non-strict policies
    and demonstrate that \brand achieves the highest level of security without sacrificing performance.
    \item \textbf{Evaluations on modern hardware:} We evaluate \brand on the latest hardware
    with nested translation support using high-end 400 Gbps networking workloads.
    \item \textbf{Open-source commitment:} We release our modified guest OS,
    experimental framework, and full implementation artifacts to support reproducibility.
\end{itemize}
% \begin{itemize}
%     % \item We identify the root cause of performance degradation in existing vIOMMU designs and analyze the security implications of different invalidation policies.
%     % \item We design \brand, a novel vIOMMU design that significantly improves performance while maintaining strict security guarantees.
%     % \item We implement \brand in the Linux kernel and evaluate its performance with various workloads, demonstrating significant improvements over existing designs.
%     \item \textbf{Overhead characterization.}
%     We provide a detailed analysis over the isnefficiencies of existing vIOMMU designs, identifying the root cause of performance degradation.
%     \item \textbf{Close-to-no-vIOMMU performance.}
%     We demonstrate that \brand achieves performance close to no-vIOMMU execution while preserving strict vIOMMU-level security guarantees.
%     \item \textbf{Strict-level security.}
%     We quantify the practical security compromises introduced by existing non-strict IOMMU invalidation policies,
%     and demonstrate that \brand can achieve the so far the best level of security without sacrificing performance.
%     \item \textbf{Evaluations on modern hardwares.}
%     We experimented \brand on most recent hardware with nested translation support and high-end 400 Gbps networking workloads.
%     % \item \textbf{Strict-equivalent security.}
%     % We show that \brand preserves the same security properties as strict invalidation while removing key performance bottlenecks.
%     \item \textbf{Open-source commitment.}
%     We release our modified guest OS, experimental framework, and full implementation artifacts to support reproducibility and community adoption.
% \end{itemize}
% For other drivers to opt in \brand, it only requires code change at the scale of a few hundred


% siyuanc3@nexus01:~/vanilla-source-code/6.12.9-iommu$ git diff main > main.patch
% siyuanc3@nexus01:~/vanilla-source-code/6.12.9-iommu$ git apply --stat main.patch
%  drivers/iommu/dma-iommu.c                       |  143 +++---
%  drivers/iommu/dma-iommu.h                       |   49 ++
%  drivers/iommu/intel/cache.c                     |  597 +++++++++++++++++++++++
%  drivers/iommu/intel/iommu.c                     |   95 ++++
%  drivers/iommu/intel/iommu.h                     |   33 +
%  drivers/iommu/iommu.c                           |    1 
%  drivers/net/ethernet/mellanox/mlx5/core/en_tx.c |  373 +++++++++++++-
%  include/linux/dma-mapping.h                     |   63 ++
%  include/linux/iommu-dma.h                       |   11 
%  include/linux/iommu.h                           |    8 
%  kernel/dma/mapping.c                            |   62 ++
%  make_and_install.sh                             |   13 +
%  12 files changed, 1319 insertions(+), 129 deletions(-)
